%%%
%%% Radical "⽎"
%%%
\section*{Radical 79: ``⽎''}\addcontentsline{toc}{section}{Radical 79: ⽎}\addcontentsline{loh}{figure}{\#\#\#\# 79: ⽎}

%%%%%%%%%% 殴 %%%%%%%%%%
\subsection*{殴}\addcontentsline{loh}{figure}{殴}

\begin{Entry}{殴}{8}{⽎}
  \begin{Phonetics}{殴}{ou1}
    \definition{v.}{espancar; bater; acertar}
  \end{Phonetics}
\end{Entry}

\begin{Entry}{殴打}{8,5}{⽎,⼿}
  \begin{Phonetics}{殴打}{ou1da3}[][HSK 7-9]
    \definition{v.}{espancar; bater; bater em alguém}
  \end{Phonetics}
\end{Entry}

%%%%%%%%%% 段 %%%%%%%%%%
\subsection*{段}\addcontentsline{loh}{figure}{段}

\begin{Entry}{段}{9}{⽎}
  \begin{Phonetics}{段}{duan4}[][HSK 2]
    \definition*{s.}{Sobrenome: Duan}
    \definition{clas.}{parte; seção; segmento; usado para dividir objetos em várias partes | passagem; parágrafo; parte de algo que tem características de continuidade | seção; período; usado para uma certa distância no tempo ou no espaço}
    \definition{s.}{nível; dan (no judô, weiqi, etc.) | seção (como nível administrativo em uma mina ou fábrica) | parte; etapa; estágio}
    \definition{v.}{cortar; separar}
  \end{Phonetics}
\end{Entry}

\begin{Entry}{段落}{9,12}{⽎,⾋}
  \begin{Phonetics}{段落}{duan4luo4}[][HSK 7-9]
    \definition[个]{s.}{seção; parágrafo; (artigo, evento) dividido em seções de acordo com o conteúdo | fase; estágio; o estágio relativamente independente de trabalho, coisas, etc.}
  \end{Phonetics}
\end{Entry}

%%%%%%%%%% 殷 %%%%%%%%%%
\subsection*{殷}\addcontentsline{loh}{figure}{殷}

\begin{Entry}{殷}{10}{⽎}
  \begin{Phonetics}{殷}{yan1}
    \definition{adj.}{vermelho"-escuro}
  \end{Phonetics}
  \begin{Phonetics}{殷}{yin1}
    \definition*{s.}{Dinastia Yin (1300--1046 a.C.; nome dado à última parte da Dinastia Shang) | Sobrenome: Yin}
    \definition{adj.}{Literário: abundante; rico; farto | Literário: ansioso; ardente; profundo | Literário: hospitaleiro; atento | grandioso; magnífico; grande | numeroso}
  \end{Phonetics}
  \begin{Phonetics}{殷}{yin3}
    \definition{s.}{Onomatopeia: estrondo do trovão}
  \end{Phonetics}
\end{Entry}

\begin{Entry}{殷勤}{10,13}{⽎,⼒}
  \begin{Phonetics}{殷勤}{yin1qin2}[][HSK 7-9]
    \definition{adj.}{solícito; ansiosamente atencioso; desejoso de agradar; caloroso e atencioso}
  \synonymref{热情}{re4qing2}
  \antonymref{怠慢}{dai4man4}
  \antonymref{冷淡}{leng3dan4}
  \end{Phonetics}
\end{Entry}

%%%%%%%%%% 殿 %%%%%%%%%%
\subsection*{殿}\addcontentsline{loh}{figure}{殿}

\begin{Entry}{殿}{13}{⽎}
  \begin{Phonetics}{殿}{dian4}
    \definition*{s.}{Sobrenome: Dian}
    \definition[座]{s.}{salão; palácio; templo}
    \definition{v.}{fechar a retaguarda (em uma marcha)}
  \end{Phonetics}
\end{Entry}

\begin{Entry}{殿堂}{13,11}{⽎,⼟}
  \begin{Phonetics}{殿堂}{dian4tang2}[][HSK 7-9]
    \definition{s.}{palácio; templo; santuário | salão do palácio; salão do templo}
  \end{Phonetics}
\end{Entry}

%%%%%%%%%% 毁 %%%%%%%%%%
\subsection*{毁}\addcontentsline{loh}{figure}{毁}

\begin{Entry}{毁}{13}{⽎}
  \begin{Phonetics}{毁}{hui3}[][HSK 6]
    \definition{v.}{destruir; arruinar; danificar | (dialeto)  transformar, remodelar um item antigo em outra coisa, geralmente roupas | queimar | difamar; caluniar}
  \end{Phonetics}
\end{Entry}

\begin{Entry}{毁灭}{13,5}{⽎,⽕}
  \begin{Phonetics}{毁灭}{hui3mie4}[][HSK 7-9]
    \definition{v.}{arruinar; destruir; exterminar; destruir ou eliminar completamente}
  \end{Phonetics}
\end{Entry}

\begin{Entry}{毁坏}{13,7}{⽎,⼟}
  \begin{Phonetics}{毁坏}{hui3huai4}[][HSK 7-9]
    \definition{v.}{danificar; devastar; degradar; destroçar; estilhaçar; vandalizar}
  \end{Phonetics}
\end{Entry}

%%%%%%%%%% 毅 %%%%%%%%%%
\subsection*{毅}\addcontentsline{loh}{figure}{毅}

\begin{Entry}{毅}{15}{⽎}
  \begin{Phonetics}{毅}{yi4}
    \definition{adj.}{firme; resoluto; inabalável}
  \end{Phonetics}
\end{Entry}

\begin{Entry}{毅力}{15,2}{⽎,⼒}
  \begin{Phonetics}{毅力}{yi4li4}[][HSK 7-9]
    \definition{s.}{coragem; vontade; resistência; força de vontade; perseverança; o espírito de nunca desistir e perseverar diante das dificuldades}
  \synonymref{坚韧}{jian1ren4}
  \synonymref{顽强}{wan2qiang2}
  \synonymref{意志}{yi4zhi4}
  \antonymref{动摇}{dong4yao2}
  \antonymref{软弱}{ruan3ruo4}
  \antonymref{颓废}{tui2fei4}
  \end{Phonetics}
\end{Entry}

\begin{Entry}{毅然}{15,12}{⽎,⽕}
  \begin{Phonetics}{毅然}{yi4ran2}[][HSK 7-9]
    \definition{adv.}{firmemente; resolutamente; com determinação}
  \seealsoref{坚决}{jian1jue2}
  \synonymref{果断}{guo3duan4}
  \synonymref{坚决}{jian1jue2}
  \antonymref{犹豫}{you2yu4}
  \end{Phonetics}
\end{Entry}

%%%%% EOF %%%%%

